% Source: https://artofproblemsolving.com/community/c3h1710118_junior_regional_olympiad__fbh_2017_grade_6_problem_1
% Stable identity: aops-topic-1710118-q01
\textbf{Useful lists.} Multiples of $7$: $7,14,\dots,98$ ($14$ numbers). Numbers containing $7$: $7,17,\dots,97$ and $70,\dots,79$, that is $10+10-1=19$ numbers, of which $7,70,77$ are multiples of $7$; the other $16$ are $17,27,37,47,57,67,71,72,73,74,75,76,78,79,87,97$. Multiples of $5$: $20$ numbers. Numbers containing $5$ but not divisible by $5$: $51,52,53,54,56,57,58,59$ ($8$ numbers).

\textbf{(a)} Faris takes $14+16=30$ cards. Of the $20$ multiples of $5$, Faris already has $35,70$ (multiples of $7$) and $75$ (contains $7$), so Lamija takes $17$. Of $51,52,53,54,56,57,58,59$, Faris has $56$ and $57$, so Lamija takes $6$ more: $23$ cards in all. \textbf{Faris has $\boxed{7}$ more cards} ($30$ against $23$).

\textbf{(b)} Lamija takes all $20+8=28$ of her cards. Of the multiples of $7$, she already has $35,56,70$, so Faris takes $11$; of the other $16$ numbers containing $7$, she has $57$ and $75$, so Faris takes $14$: $25$ cards in all. \textbf{Lamija has $\boxed{3}$ more cards} ($28$ against $25$).
